%% Thesis / Research Proposal
%% Build:  latexmk -pdf thesis_proposal.tex
%%   (or)  pdflatex -> bibtex -> pdflatex x2

\documentclass[12pt,letterpaper]{article}

%% ---------------------------------------------------------------------
%% Page geometry & spacing
%% ---------------------------------------------------------------------
\usepackage[margin=1in]{geometry}
\usepackage{setspace}
\onehalfspacing

%% ---------------------------------------------------------------------
%% Fonts & micro-typography
%% ---------------------------------------------------------------------
\usepackage[T1]{fontenc}
\usepackage[utf8]{inputenc}
\usepackage{lmodern}
\usepackage{microtype}

%% ---------------------------------------------------------------------
%% Math & logic
%% ---------------------------------------------------------------------
\usepackage{amsmath,amssymb,amsthm}
\usepackage{mathpartir}

%% ---------------------------------------------------------------------
%% Code listings
%% ---------------------------------------------------------------------
\usepackage{listings}
\lstdefinestyle{whyrel}{
  basicstyle=\small\ttfamily,
  columns=fullflexible,
  keepspaces=true,
  showstringspaces=false,
  breaklines=true,
  frame=single,
  rulecolor=\color{gray!40},
}
\lstset{style=whyrel}

%% ---------------------------------------------------------------------
%% Tables, figures, references
%% ---------------------------------------------------------------------
\usepackage{booktabs}
\usepackage{graphicx}
\usepackage{xcolor}
\usepackage[hidelinks,
            colorlinks=true,
            linkcolor=blue!60!black,
            citecolor=blue!60!black,
            urlcolor=blue!60!black]{hyperref}
\usepackage[numbers,sort&compress]{natbib}
\setcitestyle{authoryear}
%% ---------------------------------------------------------------------
%% Section formatting
%% ---------------------------------------------------------------------
\usepackage{titlesec}
\titleformat{\section}{\large\bfseries}{\thesection.}{0.5em}{}
\titleformat{\subsection}{\normalsize\bfseries}{\thesubsection.}{0.5em}{}

%% ---------------------------------------------------------------------
%% Title page helpers
%% ---------------------------------------------------------------------
\newcommand{\proposaltitle}{An Interactive Alignment Workflow for
  2-Safety and 2-Liveness Verification}
\newcommand{\proposalauthor}{Jude Kanjamala}
\newcommand{\proposaldate}{June 2026}

%% ---------------------------------------------------------------------

\begin{document}

%% =====================================================================
%% Title page
%% =====================================================================
\begin{titlepage}
  \centering
  \vspace*{2cm}

  {\LARGE\bfseries Research Proposal\par}
  \vspace{1.5cm}

  {\large\bfseries \proposaltitle \par}
  \vspace{1.5cm}

  {\large \proposalauthor \par}
  \vspace{0.5cm}

  {\normalsize Department of Computer Science \par}
  \vspace{2cm}

  % {\normalsize Submitted in partial fulfillment of the requirements\\
  %  for the degree of Master of Science in Computer Science \par}
  % \vspace{2cm}

  {\normalsize \proposaldate \par}
\end{titlepage}

%% =====================================================================
%% Abstract
%% =====================================================================
\begin{abstract}
\noindent
  % TODO. Sketch:
  %  - context: WhyRel, relational region logic, biprograms over Why3
  %  - problem: prelude/stdlib leaned on cardinality + fmap axioms and
  %    on ITP sessions bound to a pinned prover, hurting SMT robustness
  %  - contribution 1: SMT-first re-encoding (raw map.Map + explicit Fset
  %    domains; in-source `by` proof seeds replacing ITP sessions)
  %  - contribution 2: an interactive alignment workflow with push-button
  %    verification of the aligned biprogram
  %  - results: per-solver robustness and latency gains; ablation
  %    separating the encoding from the proof seeds
  This proposal is under preparation.
\end{abstract}

\tableofcontents
\newpage

%% =====================================================================
\section{Introduction and Motivation}
\label{sec:intro}
%% =====================================================================


There has been several works focused on automated alignment of programs.
\cite{naumann_isola_2020} related section gives more details on
the development of the field.

\cite{misu_2024_fse} Our main aim is to evaluate how well LLMs can synthesize
correct Dafny code, including the formal specifications and validation
conditions so that the code passes the Dafny verifier: as part of that effort,
we are also interested in how the additional logical rigor required by
verification enhances or degrades generation performance. 
- We conduct the first empirical study of LLMs synthesizing verifiable Dafny methods.
- We demonstrate that a prompt following the principles of Chain of Thought (CoT) [Wei
et al. 2022] with retrieval augmentation generated [Lewis et al. 2020] semantically similar few shot examples explaining a problem decomposition step-by-step, can synthesize
verified and correct Dafny methods with meaningful specifications for 58\% of problems
in our test dataset.
Using 178 problems from the MBPP dataset, we prompt two contemporary models
(GPT-4 and PaLM-2) to synthesize Dafny methods including annotations. We use
three different types of prompts: a direct Contextless prompt; a Signature
prompt that includes a method signature and test cases, and a Chain of Thought
(CoT) prompt that decomposes the problem into steps and includes retrieval
augmentation generated example problems and solutions. Our results show that
GPT-4 performs better than PaLM-2 on these tasks, and that both models perform
best with the retrieval augmentation generated CoT prompt. GPT-4 was able to
generate verified, human-evaluated, Dafny methods for 58\% of the problems,
however GPT-4 managed only 19\% of the problems with the Contextless prompt, and
even fewer (10\%) for the Signature prompt. 

- We release the dataset MBPP-DFY-153, a collection of 153 programming problems
with specifications, solutions and tests in/for Dafny, based on the MBPP (Mostly Basic
Python Programming) dataset curated by Google Research [Austin et al. 2021].
50 that we wrote manually and 103 synthesized by GPT-4. 

Given that we were already using LLMs to find solutions, we were also able to
employ the same LLMs to find prompts, a technique known as prompt chaining [Wei
et al. 2022] or model cascading [Dohan et al. 2022]. For each new problem, in a
first phase we prompt the LLM to generate method specifications (signature and
postconditions). Then in a second phase, we used the results of the first phase
to prompt for Dafny method implementations. The inclusion of
semantically-similar context in prompts is known as Retrieval Augmented
Generation (RAG) [Lewis et al. 2020]. Figure 3 demonstrates the resulting
high-level architecture of our RAG pipeline for prompt chaining.

in the first phase we prompt the LLM to generate the method specifications
---signature and postconditions. To teach LLM how to generate such specifications,
we provide three examples of problem descriptions, method signatures, and
postconditions, taken from our MBPP-DFY-50 dataset. We select these examples
based on the semantic similarity of the problem descriptions. Figure 4 (top)
shows an example of a first phase specification prompt. In the second phase, we
search for semantically similar Dafny methods in our MBPP-DFY-50 dataset,
incorporating the LLM's first phase response. We use the OpenAI embedding API to
compute semantic similarity based on cosine distance. For the specification
generated in the first phase, we select the top five Dafny similar examples from
our MBPP-DFY-50 dataset. We use five examples so that the whole prompt fits
inside the LLMs' context length. With these examples, we then prepare a CoT
prompt that first instructs the LLM to synthesize the specifications according
the problem decomposition steps, and then synthesize a Dafny method that
satisfies these specifications

we use verify@k metric (similar to pass@k), where k sample methods are synthesized
per problem; a problem is considered to be verified if any attempt verifies. s. In order to assess the semantic correctness of the generated verification
code, we manually reviewed all verified methods. We also manually inspected all the generated
methods that failed to verify, to identify what types of errors prevented verification.

Variation by temperature. We chose a random sample of 50 problems from our test dataset MBPP-san-DFY of
178 problems employing the sampling procedures with 90% confidence interval and 10% margin of
error. Next, for each type of prompt, we executed these 50 problems at four
different temperatures $T$ $\in$ {0.0, 0.25, 0.5, 0.75} in both GPT-4 and PaLM-2

parse errors and type errors decreased with RAG.

Jason Wei, Xuezhi Wang, Dale Schuurmans, Maarten Bosma, Brian Ichter, Fei Xia, Ed H. Chi, Quoc V. Le, and Denny Zhou.
2022. Chain-of-Thought Prompting Elicits Reasoning in Large Language Models. In Advances in Neural Information
Processing Systems 35: Annual Conference on Neural Information Processing Systems 2022, NeurIPS 2022, New Orleans, LA,
USA, November 28 - December 9, 2022, Sanmi Koyejo, S. Mohamed, A. Agarwal, Danielle Belgrave, K. Cho, and A. Oh
(Eds.). \url{http://papers.nips.cc/paper_files/paper/2022/hash/9d5609613524ecf4f15af0f7b31abca4-Abstract-Conference.html}

David Dohan, Winnie Xu, Aitor Lewkowycz, Jacob Austin, David Bieber, Raphael Gontijo Lopes, Yuhuai Wu, Henryk
Michalewski, Rif A. Saurous, Jascha Sohl-Dickstein, Kevin Murphy, and Charles Sutton. 2022. Language Model Cascades.

Patrick S. H. Lewis, Ethan Perez, Aleksandra Piktus, Fabio Petroni, Vladimir
Karpukhin, Naman Goyal, Heinrich K\"uttler, Mike Lewis, Wen-tau Yih, Tim
Rockt\"aschel, Sebastian Riedel, and Douwe Kiela. 2020. Retrieval-Augmented
Generation for Knowledge-Intensive NLP Tasks. In Advances in Neural Information
Processing Systems 33: Annual Conference on Neural Information Processing
Systems 2020, NeurIPS 2020, December 6-12, 2020, virtual, Hugo Larochelle,
Marc'Aurelio Ranzato, Raia Hadsell, Maria-Florina Balcan, and Hsuan-Tien Lin
(Eds.). \url{https://proceedings.neurips.cc/paper/2020/hash/6b493230205f780e1bc26945df7481e5-Abstract.html}

\cite{xu2025localsuccessdoescompose} DAFNYCOMP
focuses on programs composed of multiple interacting functions with necessary data dependencies,
requiring LLMs to produce specifications that ensure correctness across component boundaries.

Cite wu et all neurips 2022 autoformalization, 

\cite{chenyuan_2025_oopsla}summarize common proof strategies used by Verus
experts and is added to the prompt. AutoVerus to orchestrate these agents
following the proof-development process of Verus experts. Human experts do not aim to write a perfect proof like the one in Figure
1 in one attempt, neither does AutoVerus.

\cite{sun_2024_saiv} CLover bench as a paradigm for programming and verification that is highly amenable to llm use for synthesis of program from annotation, annotation from natural language description and inbetween translations.

\cite{loughridge2025dafnybench} focuses on annotation synthesis and thus more aligned with our own focus. It provides a benchmarks called Dafnybench having public scraped dafny programs along with cloverbench and dafny-synthesis benchmark. DafnyBench benchmark contains 782 ground-truth stand-alone Dafny programs that
compile

For hyperparameters, we set \texttt{max\_tokens} = 4096, which corresponds to the lowest
max output token limit among all the evaluated models, and we set temperature =
0.3. We gave each model up to n = 10 attempts at a given file. If it succeeded
on an attempt before the nth , it would be early stopped. If the model failed on
any of the intermediate attempts, it received the Dafny error message and was
asked to fill in the annotations again with the error message taken into
consideration. If it failed on all n attempts, it was considered to fail on that
specific test program.

innovation on the algorithmic side: out-of-the-box LLMs
provide a floor but not a ceiling for possible performance on this benchmark. For example, fine-tuning
or search-based inference-time algorithms might boost models' performances on this benchmark
(Brandfonbrener et al., 2023).

also look at \cite{mugnier_2025_oopsla}

(actually lift the dafnybench setup for ours)
- Reserach Questions?
- How much does the setup improve over teh baseline prompting.
- What is the test data (all whyrel examples) i.e benchmarks
- Task is filling annotation and finding alignments.
- pilot study.
- instead of using rag with examples, have the llms access to the internet.
- SMT solver cannot arbitrary Valid VCs and requires helping assertions.
- Prior work have demonstrated applicability of llms towards the task of synthesising intermediate assertions.
- Quantitative 
  - verified / unverified for different temperatures
- Qualitative Analysis
  - quality of preconditions, postcondition, invariants and alignments.
- types of failure
- Threats to validity
This work is the first application of llm to task of alignment selection and demonstrate the comlementary nature of llms and verifiers.

- Wanted a full loop discharge (can be put in future work section)
- for modules as well \cite{zhong_2025_LMPL}


The remainder of this proposal is organized as follows.
Section~\ref{sec:background} provides background on relational verification
and the WhyRel toolchain.
Section~\ref{sec:proposed} describes the proposed approach and research
questions.
Section~\ref{sec:prelim} summarizes preliminary results.
Section~\ref{sec:related} surveys related work.
Section~\ref{sec:plan} outlines the research plan and timeline.

%% =====================================================================
\section{Background}
\label{sec:background}
%% =====================================================================

\subsection{Relational Region Logic and WhyRel}

% TODO: expand background on RRL, biprograms, Why3 encoding.

\subsection{SMT Encoding Considerations}
\label{sec:encoding}

% raw map.Map + explicit Fset domains; portfolio (alt-ergo/z3/cvc5);
% in-source `by` extensionality and witness seeds.

%% =====================================================================
\section{Proposed Work}
\label{sec:proposed}
%% =====================================================================

\subsection{Objectives}

\begin{enumerate}
  \item Automate alignment construction and annotation suggestion to reduce
        manual proof effort without sacrificing verifiability.
  \item Evaluate whether an SMT-first re-encoding of the WhyRel prelude
        improves solver robustness and latency across a representative
        benchmark suite.
  \item Characterize the interaction between the re-encoding and automated
        alignment, measuring whether the two contributions compose favorably.
\end{enumerate}

\subsection{Interactive Alignment Workflow}
\label{sec:align}



% TODO: brief description of the implementation.

%% =====================================================================
\section{Preliminary Results}
\label{sec:prelim}
%% =====================================================================

% TODO: suite-wide before/after; tiling latency ablation (encoding vs seeds).

%% =====================================================================
\section{Related Work}
\label{sec:related}
%% =====================================================================

While the current workflow is focused on bimethods, there has been work on
whole-project synthesis~\cite{zhong_2025_LMPL}.

%% =====================================================================
\section{Research Plan and Timeline}
\label{sec:plan}
%% =====================================================================

% TODO: milestones and expected completion dates.

%% =====================================================================
\section{Conclusion}
\label{sec:conclusion}
%% =====================================================================

%% =====================================================================
%% Bibliography
%% =====================================================================
\bibliographystyle{plainnat}
\bibliography{refs}

\end{document}
